% Incorporación ampliativa: CG-006, SR-032--SR-035 del inventario REV05.
% Residencia: inmediatamente después de la construcción de R36, antes de G9.
% Propietario computacional: 16_CIERRE_GLOBAL_HMT_MD_2026-07-22/
% 01_SELECTOR_CONSTANTES/selector_interno_constantes.py, líneas 31--63,
% 131--180, 185--240 y 267--321.
% Recuperación y cambio de representación: output/REVISION_EDITORIAL_HMT_MD_
% 20260905/03_PRUEBAS/verificar_generacion_infinita_nonadica.py,
% líneas 54--77, 651--723 y 974--1147.
% Exposición de origen: ampliaciones_sucesoras_20260824/parte_i_ii/owners/
% U008_orbitas_elevacion_r36.tex, líneas 210--358.
% Estatuto: RESULTADO_RECUPERADO; exposición de enumeración y naturalidad
% reunida sin promover la concordancia en R36 a equivalencia global.

\section[Selección y transporte de la primera frontera]{Selección combinatoria y transporte de representación de la primera frontera conjunta}
\label{r36c:sec:frontera}

La primera frontera conjunta recibe los bloques ya producidos por las
semillas APP, la reducción trítica y el calendario de elevación del TPK.
Su construcción conserva la diferencia entre tres operaciones: determinar
las distribuciones admisibles, orientar sus dos canales laterales y
transportar la incidencia al sistema de coordenadas excepcional. La
presente exposición desarrolla estas operaciones sobre el mismo estado
finito. Los nombres de las coordenadas reales designan sus lectores
posteriores; las condiciones de selección que siguen actúan sobre
palabras ternarias, matrices de transporte e incidencias enteras.

\subsection{Estado regional anterior a la frontera}

Sea $\mathcal W\subset\mathbb F_3^6$ el repertorio de $243$ palabras
visibles obtenido del censo de emisiones APP--TRIT--TPK. Para un vector
fila $s\in\mathcal W$, definimos
\[
 b_0(s)=s,\quad b_1(s)=sL_0,\quad b_2(s)=sL_0^2,\quad
 b_3(s)=sL_0^2L_1,\quad b_4(s)=sL_0^2L_1^2,
\]
con operaciones en $\mathbb F_3$. Las matrices previamente derivadas de
los registros regionales son
\[
L_0=\begin{pmatrix}
2&2&2&1&2&1\\2&1&2&2&1&1\\1&0&0&1&0&2\\
0&0&1&1&1&0\\2&2&2&0&2&1\\2&1&2&1&0&0
\end{pmatrix},\qquad
L_1=\begin{pmatrix}
2&1&2&0&2&2\\1&2&1&1&2&0\\1&2&1&1&1&2\\
0&1&0&0&2&1\\2&0&2&1&0&1\\0&2&2&2&1&1
\end{pmatrix}.
\]
El selector regional anterior entrega
\[
s_{\rm cl}=010211,\qquad s_{\rm pr}=201101,
\]
mediante las condiciones
$b_2(s_{\rm cl})=b_0(s_{\rm cl})$,
$b_4(s_{\rm pr})=b_2(s_{\rm pr})$ y
$b_1(s_{\rm cl})=b_3(s_{\rm pr})$.
En esta etapa, los bloques corrientes son
\[
b_4(s_{\rm cl})=110221,\qquad
b_4(s_{\rm pr})=102011.
\]
La concatenación conserva los cinco bloques de cada registro; el vector
$b_4$ es la entrada de seis coordenadas que recibe la relación de frontera.

Para expresar fielmente la relación recuperada se utiliza la matriz de
incidencia en su representación de origen,
\begin{equation}
A^{\rm o}=\begin{pmatrix}
0&1&1&1&1&1\\1&0&1&1&2&2\\1&1&0&2&1&2\\
1&1&2&0&2&1\\1&2&1&2&0&1\\1&2&2&1&1&0
\end{pmatrix}.
\label{r36c:eq:ao}
\end{equation}
Se verifica $(A^{\rm o})^2=-I$ en $M_6(\mathbb F_3)$; por tanto,
$(A^{\rm o})^{-1}=-A^{\rm o}$. Más adelante se transportará esta
representación, junto con sus vectores, a la utilizada por el lector
excepcional.

\subsection{Ecuación de compatibilidad y saturación columnar}

El eje de la frontera queda fijado por
$z=b_4(s_{\rm pr})=102011$. Para cada $s\in\mathcal W$ se define
\begin{align}
m(s)&=b_4(s)A^{\rm o}L_1^3,\\
q(s)&=m(s)+z,\\
\widetilde q(s)&=s(A^{\rm o})^{-1}L_0^3(A^{\rm o})^3.
\label{r36c:eq:compat-datos}
\end{align}
La condición de compatibilidad es $q(s)=\widetilde q(s)$. Las filas
laterales $x,y\in\mathbb F_3^6$ satisfacen $x+y=m(s)$; la frontera
ordenada es $B=(x,y,z)^{\mathsf T}$.

Para formar el residuo y el acarreo se usan los representantes enteros
$0,1,2$ de cada trit. Si $S_j=x_j+y_j+z_j$, se exige
\[
q_j=S_j\bmod3,\qquad c_j=\frac{S_j-q_j}{3}=1,
\]
y la ocupación columnar viene dada por
\begin{equation}
\mathbf1_{\{x_j\ne0\}}+\mathbf1_{\{y_j\ne0\}}+\mathbf1_{\{z_j\ne0\}}
=2+\mathbf1_{\{z_j\ne0,\ m_j(s)\ne2\}}.
\label{r36c:eq:ocupacion}
\end{equation}
La neutralidad dual completa la condición mediante
\begin{equation}
BA^{\rm o}\mathbf1=0\quad\hbox{en }\mathbb F_3^3.
\label{r36c:eq:neutralidad}
\end{equation}
El residuo, el acarreo y la ocupación intervienen conjuntamente: la
reducción modular conserva el residuo y el registro entero distingue las
distribuciones que lo producen.

\begin{lemma}[Resolución de la compatibilidad regional]
\label{r36c:lem:compatibilidad}
La ecuación $q(s)=\widetilde q(s)$ tiene, en $\mathbb F_3^6$, las tres
soluciones
\[
020010,\qquad121200,\qquad222120.
\]
Su intersección con el repertorio visible es
$\{121200,222120\}$.
\end{lemma}
\begin{proof}
La ecuación se escribe $sD=z$, donde
\[
D=(A^{\rm o})^{-1}L_0^3(A^{\rm o})^3
   -L_0^2L_1^2A^{\rm o}L_1^3
 =\begin{pmatrix}
2&1&1&2&2&2\\0&2&0&0&0&1\\0&0&0&0&1&0\\
1&2&2&2&2&0\\1&2&2&0&1&2\\1&2&0&2&2&2
\end{pmatrix}.
\]
La eliminación sobre $\mathbb F_3$ da la parametrización
\[
s=(t,2,t,2t,1+2t,0),\qquad t\in\mathbb F_3.
\]
Su sustitución verifica las seis ecuaciones y agota sus soluciones. En
el repertorio generado $\mathcal W$, las dos palabras correspondientes a
$t=1,2$ están presentes y la correspondiente a $t=0$ está ausente. Esta
última comprobación es una consulta exacta al censo de palabras visibles,
anterior a cualquier lectura arquimediana.
\end{proof}

\begin{proposition}[Determinación del eje de autoescala]
La compatibilidad, el acarreo unitario y la ocupación columnar seleccionan
$s_{\rm au}=121200$ entre las dos soluciones visibles del
lema~\ref{r36c:lem:compatibilidad}. Se obtiene
\[
b_4(s_{\rm au})=010200,\qquad
m(s_{\rm au})=210112,\qquad q=012120,
\]
con ocupación $223232$.
\end{proposition}
\begin{proof}
Para $s=222120$ se calculan
$b_4(s)=222000$, $m(s)=201010$ y $q(s)=000021$.
En su tercera columna, $m_3=1$ y $z_3=2$. La ocupación requerida es tres,
de modo que $x_3,y_3$ deben ser ambos distintos de cero. La única pareja
con $x_3+y_3\equiv1\pmod3$ es $(2,2)$; su suma entera con $z_3=2$ es
seis y produce acarreo dos. Contradice $c_3=1$.

Para $s=121200$ las seis columnas admiten, respectivamente, las parejas
\[
\begin{array}{c|cccccc}
j&1&2&3&4&5&6\\\hline
(x_j,y_j)&(0,2),(2,0)&(2,2)&(1,2),(2,1)&(2,2)&(2,2)&(0,2),(2,0).
\end{array}
\]
Cada pareja satisface las ecuaciones de residuo, acarreo y ocupación.
Hay, por tanto, ocho disposiciones antes de aplicar la neutralidad dual.
\end{proof}

\begin{theorem}[Enumeración exacta de las orientaciones de frontera]
\label{r36c:thm:dos}
Las condiciones anteriores determinan exactamente
\[
B_- =\begin{pmatrix}021222\\222220\\102011\end{pmatrix},
\qquad
B_+ =\begin{pmatrix}222220\\021222\\102011\end{pmatrix}.
\]
Ambas tienen residuo $012120$, acarreo $111111$ y ocupación $223232$.
\end{theorem}
\begin{proof}
La prueba anterior reduce la enumeración a ocho posibilidades. Puesto que
$A^{\rm o}\mathbf1=(2,1,1,1,1,1)^{\mathsf T}$, la neutralidad de una
fila $v$ equivale a $2v_1+v_2+\cdots+v_6=0$ en $\mathbb F_3$.
Su evaluación en las ocho disposiciones es
\[
\begin{array}{c|c|c}
x&y&(xA^{\rm o}\mathbf1,yA^{\rm o}\mathbf1,zA^{\rm o}\mathbf1)\\\hline
021220&222222&(1,2,0)\\
021222&222220&(0,0,0)\\
022220&221222&(2,1,0)\\
022222&221220&(1,2,0)\\
221220&022222&(2,1,0)\\
221222&022220&(1,2,0)\\
222220&021222&(0,0,0)\\
222222&021220&(2,1,0).
\end{array}
\]
Las filas segunda y séptima son las únicas soluciones. Las sumas
columnarias enteras de cualquiera de ellas son $(3,4,5,4,5,3)$, que
recuperan el residuo y el acarreo declarados.
\end{proof}

\subsection{Longitud dirigida y orientación de la frontera}

El transporte regional define un grafo dirigido sobre $\mathbb F_3^6$,
con aristas $v\to vL_0$ y $v\to vL_1$. Para vectores $u,v$ alcanzables
entre sí, sea
\[
d(u,v)=\min\{n\ge0:\exists\epsilon_1,\ldots,\epsilon_n\in\{0,1\},
\ v=uL_{\epsilon_1}\cdots L_{\epsilon_n}\}.
\]
Esta longitud es un dato combinatorio del transporte. Su suma sobre las
tres filas conserva las etiquetas regionales:
\[
\mathcal L(B)=d(110221,B_1)+d(102011,B_2)+d(010200,B_3).
\]

\begin{proposition}[Selección por longitud dirigida]
\label{r36c:prop:minimo}
Se cumple $\mathcal L(B_-)=23$ y $\mathcal L(B_+)=20$. En el dominio
de las dos fronteras admisibles, el mínimo es único y selecciona $B_+$.
\end{proposition}
\begin{proof}
Los caminos mínimos y sus longitudes son
\[
\begin{array}{c|c|c|c|c}
 &\text{origen}&\text{destino}&\text{palabra de transporte}&d\\\hline
B_-&110221&021222&01100011&8\\
 &102011&222220&10111111&8\\
 &010200&102011&1010101&7\\\hline
B_+&110221&222220&1111&4\\
 &102011&021222&001100101&9\\
 &010200&102011&1010101&7.
\end{array}
\]
Cada palabra se verifica por multiplicación sucesiva de sus matrices.
Para comprobar también la minimalidad se forman las capas disjuntas
\[
F_0(u)=\{u\},\qquad
F_{n+1}(u)=\bigl(F_n(u)L_0\cup F_n(u)L_1\bigr)
\setminus\bigcup_{j=0}^{n}F_j(u).
\]
Una inducción en $n$ demuestra que $F_n(u)$ es exactamente el conjunto
de vértices a distancia $n$: toda palabra mínima de longitud $n+1$
procede de una palabra mínima de longitud $n$, y el cociente por las
capas anteriores elimina los recorridos de menor longitud.
La enumeración obtenida con las dos matrices expuestas arriba da
\[
\begin{array}{c|rrrrrrrrrr}
u\backslash n&0&1&2&3&4&5&6&7&8&9\\\hline
110221&1&2&4&8&16&32&60&105&156&177\\
102011&1&2&3&5&10&20&37&71&119&171\\
010200&1&2&4&8&15&30&56&100&148&172.
\end{array}
\]
Los destinos de la primera tabla pertenecen por primera vez a las capas
allí indicadas. Esta enumeración es finita, utiliza únicamente los
productos módulo tres y coincide con el recorrido por anchura del
certificado ejecutable. Sumando las tres longitudes resulta $8+8+7=23$
y $4+9+7=20$.
\end{proof}

La orientación de la representación APP mediante cilindros semiabiertos
selecciona previamente la misma matriz $B_+$ en esta frontera. El
teorema~\ref{r36c:thm:dos} y la
proposición~\ref{r36c:prop:minimo} establecen así la concordancia de
ambas selecciones sobre el conjunto concreto $\{B_-,B_+\}$. El dominio
de esta concordancia es la primera frontera conjunta. La prolongación
coinductiva conserva sus propias relaciones de admisibilidad, memoria,
supervivencia y orientación en cada profundidad.

\subsection{Naturalidad del cambio de representación excepcional}

La matriz empleada por el lector excepcional en la representación activa es
\[
A^{\rm a}=\begin{pmatrix}
0&1&1&1&1&1\\1&0&1&2&2&1\\1&1&0&1&2&2\\
1&2&1&0&1&2\\1&2&2&1&0&1\\1&1&2&2&1&0
\end{pmatrix}.
\]
La permutación de coordenadas
\[
R_p(v_1,v_2,v_3,v_4,v_5,v_6)
=(v_1,v_2,v_3,v_5,v_6,v_4)
\]
transporta simultáneamente la palabra y la matriz de incidencia.

% REMATE_LOCAL_R36
\Needspace{36\baselineskip}
\begin{theorem}[Entrelazamiento de la incidencia de frontera]
\label{r36c:thm:cambio}
Para todo $v\in\mathbb F_3^6$,
\[
R_p(vA^{\rm o})=R_p(v)A^{\rm a}.
\]
La transformación es invertible y conserva la neutralidad dual, los
residuos, los acarreos y las ocupaciones, expresados en las coordenadas
transportadas. En particular, la terna seleccionada tiene las dos
representaciones
\[
\begin{array}{c|c|c}
&\text{representación de origen}&\text{representación excepcional}\\\hline
\text{clausura}&222220&222202\\
\text{propagación}&021222&021222\\
\text{autoescala}&102011&102110.
\end{array}
\]
\end{theorem}
\begin{proof}
Si $p=(1,2,3,5,6,4)$, la comparación de las entradas de ambas matrices
da $A^{\rm a}_{ij}=A^{\rm o}_{p(i),p(j)}$. Para la coordenada $j$,
\[
\bigl(R_p(v)A^{\rm a}\bigr)_j
=\sum_i v_{p(i)}A^{\rm o}_{p(i),p(j)}
=\sum_kv_kA^{\rm o}_{k,p(j)}
=\bigl(R_p(vA^{\rm o})\bigr)_j.
\]
La inversa es
$R_p^{-1}(w)=(w_1,w_2,w_3,w_6,w_4,w_5)$.
Como $R_p$ permuta coordenadas, conserva su suma y transporta
coordenada a coordenada las sumas enteras, sus residuos módulo tres,
sus cocientes y las cantidades de entradas distintas de cero. De ello
resulta la conservación de todos los registros indicados. La tabla se
obtiene aplicando $R_p$ a las tres palabras de $B_+$ y la inversa
recupera exactamente sus palabras originales.
\end{proof}

El lector posicional y el lector de incidencia permanecen asociados al
mismo estado enriquecido mediante este cambio explícito de
representación. Si $\ell_{\rm pos}$ es el lector en la representación
de origen, su expresión sobre la representación transportada es
$\ell_{\rm pos}\circ R_p^{-1}$. Esta composición conserva el orden
posicional: cambiar la representación de la incidencia y cambiar el
lector son operaciones coordinadas. La primera frontera entra así en
la prolongación con sus tres palabras, su orientación regional, su
acarreo entero y su marco excepcional recuperables.
